% !TEX root = ../Abschlussarbeit_Jan_Händl.tex
\section{Ausblick}
Die Ausblicke und möglichen Weiterentwicklungsansätze sind in Bezug auf die Microfrontend-Architektur, als auch das Domänenmodell vielfältig.
An erster Stelle steht die vollständige Überführung des IFC-Standards in GraphQL Schemata. Wie bereits in Absatz \ref{sec:graphql} erörtert, gestaltet sich aufgrund der Komplexität des IFC-Standards
eine 1:1 Abbildung des IFC-Standards in GraphQL Schemata als äußerst aufwendig. Hier könnte entweder eine Abstraktion des IFC-Standards zielführend sein \cite{wolf_20206_graphql}, oder die Nutzung von
größeren Sprachmodellen zur automatisierten Überführung in GraphQL Schemata. Eng damit verbunden ist die Entwicklung eines leistungsfähigen IFC-Parsers. Ein solcher Parser ist notwendig, um IFC-Dateien effizient in die Graphdatenbank importieren zu können.
Zwar existieren in der Praxis bereits verschiedene Open-Source-Parser \cite{wolf:2025:ifc2graphql} zur Verarbeitung von IFC-Dateien, allerdings ist keiner dieser Ansätze auf die spezifischen Anforderungen dieses Forschungsthemas optimiert. Es bedürfe daher eine individuelle Lösung.
Auch in Bezug auf Multi-Tenancy ergeben sich weitere Forschungsansätze. Innerhalb dieser Arbeit wurde lediglich eine clientseitige Umsetzung
innerhalb der Microfrontend-Architektur erwähnt, während GraphQL mittels seiner @auth Direktive ein Rechtemanagement auf dem Server ermöglichen könnte. Dabei könnte der Nutzen und eine mögliche Umsetzung eines Rechtemanagements innerhalb
einer Graphdatenbank thematisiert werden, oder das Zusammenspiel mit einer relationalen Datenbank. Dabei könnte auch die Erstellung eines Rechtemanagements an sich, mit notwendigen Rollen und Berechtigungen im Kontext des Inbetriebnahmemanagements, als Forschungsthema dienen.
Neben diesen möglichen Ansätzen beinhaltet auch die Microfrontend-Architektur offene Fragen, wie beispielsweise eine durchgehende \ac{CI/CD}-Pipeline zur automatisierten Bereitstellung der einzelnen Microfrontends.
Hier könnten insbesondere die Herausforderungen des Bereitstellens und Komposition der einzelnen Microfrontends im Vordergrund stehen, sowie eine mögliche automatische Skalierung, um Lastspitzen auszugleichen. Analog zu klassischen Microservices könnte dies das Potential dieser verteilten Systeme weiter ausschöpfen. \\
Als weiterer wichtiger Punkt wären mögliche Implementierungen anderer Frontendtechnologien neben React zu nennen. Da Module Federation unabhängig von Frameworks ist, könnten eine Implementierung anderer Frontendtechnologien die Technologieunabhängigkeit weiter untermauern. Im Zuge dessen wäre auch die frameworkübergreifende Interaktion zwischen den Microfrontends, in Bezug auf Kommunikation, Routing und Isolation, von Interesse. \newline
Den größten Erkenntnisgewinn verspricht jedoch die Validierung des Gesamtsystems unter realitätsnahen Bedingungen. Ein empirisches Testen der Graphdatenbank mit großen IFC-Datenmengen sowie die gleichzeitige Nutzung der Microfrontends durch eine Vielzahl von Akteuren würde den größten Informationsgehalt hinsichtlich der tatsächlichen Performance, Skalierbarkeit, Effizienz und Systemstabilität liefern. Erst durch ein solches Last- und Stressszenario lässt sich die Praxistauglichkeit der entwickelten Microfrontend-Architektur in Verbindung mit dem graphbasierten Domänenmodell nachweisen.